\chapter{Relación con los Objetivos de Desarrollo Sostenible}
\label{anx:ods}

Grado de relación del trabajo con los Objetivos de Desarrollo Sostenible de la
Agenda 2030.

\begin{table}[htbp]
    \centering
    \caption{Grado de relación del trabajo con los Objetivos de Desarrollo
             Sostenible.}
    \label{tab:ods}
    \begin{tabular}{lcccc}
        \toprule
        \textbf{Objetivos de Desarrollo Sostenibles} &
        \textbf{Alto} & \textbf{Medio} & \textbf{Bajo} &
        \textbf{No procede} \\
        \midrule
        ODS 1. Fin de la pobreza.                          &   &   &   & X \\
        ODS 2. Hambre cero.                                &   &   &   & X \\
        ODS 3. Salud y bienestar.                          &   &   &   & X \\
        ODS 4. Educación de calidad.                       &   & X &   &   \\
        ODS 5. Igualdad de género.                         &   &   &   & X \\
        ODS 6. Agua limpia y saneamiento.                  &   &   &   & X \\
        ODS 7. Energía asequible y no contaminante.        &   &   &   & X \\
        ODS 8. Trabajo decente y crecimiento económico.    &   &   &   & X \\
        ODS 9. Industria, innovación e infraestructuras.   & X &   &   &   \\
        ODS 10. Reducción de las desigualdades.            &   &   &   & X \\
        ODS 11. Ciudades y comunidades sostenibles.        &   &   &   & X \\
        ODS 12. Producción y consumo responsables.         &   & X &   &   \\
        ODS 13. Acción por el clima.                       &   &   &   & X \\
        ODS 14. Vida submarina.                            &   &   &   & X \\
        ODS 15. Vida de ecosistemas terrestres.            &   &   &   & X \\
        ODS 16. Paz, justicia e instituciones sólidas.     &   &   &   & X \\
        ODS 17. Alianzas para lograr objetivos.            &   &   &   & X \\
        \bottomrule
    \end{tabular}
\end{table}

\section*{Descripción de la alineación del trabajo con los ODS con un grado de
relación más alto}

\textbf{ODS 9. Industria, innovación e infraestructuras.} Es el objetivo con el
que este trabajo se relaciona de forma directa, y en concreto con su meta de
apoyar la investigación científica y mejorar la capacidad tecnológica. El
trabajo no aplica una técnica ya resuelta, sino que lleva un modelo generativo
entrenado en un ordenador hasta un instrumento embebido que funciona en tiempo
real y sin ordenador conectado, sobre placas de desarrollo de bajo coste. La
aportación no está en la interpolación neuronal de tablas de onda, que ya
existía en la literatura, sino en resolver los problemas que aparecen al
llevarla a un aparato: repartir el trabajo entre tres procesadores según la
escala de tiempo de cada tarea, precalcular en el ordenador lo que no cabe
ejecutar a bordo y garantizar que la información llegue íntegra entre placas.
Ese tipo de trabajo —adaptar una técnica de investigación a hardware asequible—
es justamente lo que el objetivo entiende por mejorar la capacidad tecnológica,
porque baja la barrera de entrada de una técnica que hasta ahora exigía un
ordenador dedicado.

\textbf{ODS 4. Educación de calidad.} El desarrollo completo está publicado en
un repositorio de acceso libre, y no sólo el resultado. Incluye los guiones de
preprocesado del conjunto de datos, el entrenamiento del modelo, el horneado de
la rejilla y las validaciones numéricas, además del firmware de las tres placas
y los ficheros paramétricos de la carcasa. Eso permite reproducir el trabajo
entero o reutilizar cualquiera de sus partes por separado. A ello se suma que la
memoria documenta también los caminos que no funcionaron y por qué, que suele
ser la información que más cuesta encontrar y la más útil para quien intente algo
parecido.

\textbf{ODS 12. Producción y consumo responsables.} Varias decisiones de diseño
del capítulo~\ref{cap:hardware} apuntan en esta dirección, y ninguna se tomó
pensando en la sostenibilidad, sino por razones prácticas que resultan coincidir
con ella. Las dos placas de desarrollo caras van sobre zócalos y no soldadas,
así que se recuperan intactas si la placa que las sostiene falla o si el
proyecto se desmonta. Todas las uniones de la carcasa aprietan contra una tuerca
metálica embutida en lugar de roscar en el plástico, precisamente porque el
plástico aguanta pocos montajes y la pieza dejaría de ser desmontable. El
instrumento se construye además con componentes de catálogo, sin ningún circuito
integrado a medida, de modo que cualquier pieza averiada se sustituye por otra
equivalente. Y el diseño de la carcasa es paramétrico y está publicado, así que
una pieza rota se vuelve a fabricar sin rehacer el modelo.
